% !TEX program = xelatex
% ============================================================
%  Arabic Article / Letter Template — Nibras Center
%  Compile with: xelatex main.tex (run twice for TOC)
%  Based on the ASPU_Proposal_AR project setup.
% ============================================================
%  This is the PLACEHOLDER version — it includes sample text
%  and comments explaining each section. For an empty skeleton,
%  use main_skeleton.tex instead.
% ============================================================

\documentclass[11pt,a4paper]{article}

% ---- Geometry ----
\usepackage[a4paper,margin=2.2cm,top=2.4cm,bottom=2.4cm,headheight=15pt]{geometry}

% ---- Core packages (order matters for bidi) ----
\usepackage{url}
\usepackage{amsmath}
\usepackage{graphicx}
\usepackage{array}
\usepackage{booktabs}
\usepackage{tabularx}
\usepackage{multirow}
\usepackage{enumitem}
\usepackage{float}
\usepackage{caption}
\usepackage{parskip}
\usepackage{titlesec}
\usepackage{fancyhdr}
\usepackage{setspace}

% ---- RTL / Arabic language support (must load before hyperref) ----
\usepackage{bidi}
\usepackage[no-math]{fontspec}
\usepackage[quiet]{polyglossia}

% ---- Fonts ----
% Noto Kufi Arabic is a geometric Kufi-style Arabic font from Google Noto Fonts.
% Replace "NotoKufiArabic-Regular.ttf" with your preferred font if needed.
\setmainfont[Scale=1]{NotoKufiArabic-Regular.ttf}
\newfontfamily\arabicfont[Script=Arabic,Scale=0.95]{NotoKufiArabic-Regular.ttf}
\newfontfamily\englishfont[Scale=0.95]{TeX Gyre Termes}
\setmonofont{NotoKufiArabic-Regular.ttf}

% ---- Languages ----
\setdefaultlanguage[calendar=gregorian,hijricorrection=1]{arabic}
\setotherlanguage[variant=british]{english}

% ---- Hyperref (load after polyglossia/bidi) ----
\usepackage[breaklinks,hidelinks]{hyperref}
\hypersetup{
  pdftitle={عنوان المستند},
  pdfauthor={د. اسم المؤلف}
}

% ---- Captions in Arabic ----
\addto\captionsarabic{%
  \renewcommand{\contentsname}{الفهرس}%
  \renewcommand{\figurename}{الشكل}%
  \renewcommand{\tablename}{الجدول}%
  \renewcommand{\refname}{المراجع}%
}

% ---- Section formatting ----
\titleformat{\section}{\Large\bfseries}{\thesection}{0.7em}{}
\titleformat{\subsection}{\large\bfseries}{\thesubsection}{0.6em}{}
\titleformat{\subsubsection}{\normalsize\bfseries}{\thesubsubsection}{0.5em}{}
\titlespacing*{\section}{0pt}{1.4ex plus 0.4ex}{0.9ex plus 0.3ex}
\titlespacing*{\subsection}{0pt}{1.0ex plus 0.3ex}{0.7ex plus 0.2ex}

% ---- Headers / footers ----
\pagestyle{fancy}
\fancyhf{}
\fancyhead[R]{\small\itshape عنوان مختصر للمستند}
\fancyhead[L]{\small د. اسم المؤلف}
\fancyfoot[C]{\small\thepage}
\renewcommand{\headrulewidth}{0.3pt}
\renewcommand{\footrulewidth}{0pt}

% ---- List spacing ----
\setlist[itemize]{topsep=0.3em, itemsep=0.15em, parsep=0pt}
\setlist[enumerate]{topsep=0.3em, itemsep=0.15em, parsep=0pt}

% ---- Table column types ----
\newcolumntype{L}[1]{>{\raggedright\arraybackslash}p{#1}}
\newcolumntype{C}[1]{>{\centering\arraybackslash}p{#1}}

% ---- Line spacing ----
\setstretch{1.25}

\begin{document}

% ============================================================
% TITLE PAGE
% ============================================================
\begin{titlepage}
\centering
\vspace*{2cm}
{\Huge\bfseries عنوان المستند\\[0.3em]
عنوان فرعي\par}
\vspace{1.5cm}
{\large
\begin{tabular}{rl}
\textbf{مُقدَّم إلى:} & الجهة المُستلمة\\
\textbf{مُقدَّم من:} & د. اسم المؤلف\\
\textbf{المجال:} & المجال أو الموضوع\\
\textbf{التاريخ:} & شهر سنة\\
\end{tabular}\par}
\vspace{1.5cm}
{\large\itshape \url{https://example.com}\par}
\vfill
{\small ملف مرافق: \texttt{file.md} — وصف مختصر}
\end{titlepage}

\newpage
\setcounter{page}{2}

% ============================================================
% COVER LETTER (optional)
% ============================================================
\section*{خطاب التقديم}
\addcontentsline{toc}{section}{خطاب التقديم}

\begin{flushright}
\textbf{د. اسم المؤلف}
\end{flushright}

\vspace{0.5em}
\textbf{إلى:} اسم الجهة المُستلمة

\vspace{0.5em}
\textbf{الموضوع:} موضوع الخطاب

\vspace{0.8em}
السيد المحترم،

% TODO: اكتب نص الخطاب هنا. استخدم أسلوباً رسمياً.

\vspace{1em}
وتفضلوا بقبول فائق التقدير،\\[0.5em]
\textbf{د. اسم المؤلف}\\
\url{https://example.com}

\newpage

% ============================================================
% SECTION 1
% ============================================================
\section*{عنوان القسم الأول}
\addcontentsline{toc}{section}{عنوان القسم الأول}

% TODO: اكتب محتوى القسم الأول هنا.

\subsection*{1.1 عنوان القسم الفرعي}

% TODO: اكتب المحتوى هنا.

\subsection*{1.2 عنوان القسم الفرعي}

% TODO: اكتب المحتوى هنا.

\newpage

% ============================================================
% SECTION 2
% ============================================================
\section*{عنوان القسم الثاني}
\addcontentsline{toc}{section}{عنوان القسم الثاني}

% TODO: اكتب محتوى القسم الثاني هنا.

\subsection*{2.1 عنوان القسم الفرعي}

% TODO: اكتب المحتوى هنا.

\newpage

% ============================================================
% APPENDICES
% ============================================================
\section*{الملحق الأول — عنوان الملحق}
\addcontentsline{toc}{section}{الملحق الأول}

% TODO: اكتب محتوى الملحق الأول هنا.

\section*{الملحق الثاني — عنوان الملحق}
\addcontentsline{toc}{section}{الملحق الثاني}

% TODO: اكتب محتوى الملحق الثاني هنا.

\end{document}
